% GENERATED FILE. Edit canonical lessons, then run npm run generate:books.
% canonical-source-hash: fnv1a64:c915516cc0affa08
% canonical-lessons: TE-C283-charj-ceyu, TE-C283-tolagincu, TE-C283-gurtupattu, TE-C283-vyatirekincu, TE-C283-vadincu

\chapter{To charge, To remove, To recognize, To oppose, To argue}
\label{ch:te-to-charge-to-remove-to-recognize-to-oppose-to-argue}
\hlchaptermodality{283}

\begin{chapteropening}
\textbf{By the end of this chapter:} \emph{I can say to charge, to remove, to recognize, to oppose and to argue.}

Complete the last lesson of chapter 283: I can say to charge, to remove, to recognize, to oppose and to argue.
\end{chapteropening}

\section[\texorpdfstring{\te{ఛార్జ్} \te{చేయు} (chārj cēyu)}{ఛార్జ్ చేయు (chārj cēyu)}]{\te{ఛార్జ్} \te{చేయు} (chārj cēyu) — to charge (a battery)}

\label{lesson:TE-C283-charj-ceyu}



\begin{quote}
Before the new one: say the Telugu for to share out, then the Telugu for to visit.
\end{quote}

\subsection*{You'll want to know: \te{ఛార్జ్} \te{చేయు}}
\textbf{\te{ఛార్జ్} \te{చేయు}} — \emph{chārj cēyu} — \textquotedblleft{}to charge (a battery)\textquotedblright{}.

\begin{grammarlens}[title={What you've built}]
One more word for this chapter.
\end{grammarlens}

\subsection*{Your turn}
\begin{itemize}
\raggedright
  \item \emph{Say it:} \emph{chārj cēyu}
  \item \emph{Say it:} \emph{chārj cēyu}, once more
  \item \emph{Say it:} say \emph{chārj cēyu}
  \item \emph{Recall:} say the Telugu for to share out, then the Telugu for to visit, then say \emph{chārj cēyu} again
\end{itemize}

\begin{tcolorbox}[breakable,colback=teal!4,colframe=teal!35!black,title={Before you move on}]
What does \te{ఛార్జ్} \te{చేయు} mean? (\textquotedblleft{}To charge (a battery)\textquotedblright{}.) Say it once more.
\end{tcolorbox}
\section[\texorpdfstring{\te{తొలగించు} (tolagiñcu)}{తొలగించు (tolagiñcu)}]{\te{తొలగించు} (tolagiñcu) — to remove, to delete}

\label{lesson:TE-C283-tolagincu}



\begin{quote}
Before the new one: say the Telugu for to visit, then the Telugu for to charge.
\end{quote}

\subsection*{You'll want to know: \te{తొలగించు}}
\textbf{\te{తొలగించు}} — \emph{tolagiñcu} — \textquotedblleft{}to remove, to delete\textquotedblright{}.

\begin{grammarlens}[title={What you've built}]
One more word for this chapter.
\end{grammarlens}

\subsection*{Your turn}
\begin{itemize}
\raggedright
  \item \emph{Say it:} \emph{tolagiñcu}
  \item \emph{Say it:} \emph{tolagiñcu}, once more
  \item \emph{Say it:} say \emph{tolagiñcu}
  \item \emph{Recall:} say the Telugu for to visit, then the Telugu for to charge, then say \emph{tolagiñcu} again
\end{itemize}

\begin{tcolorbox}[breakable,colback=teal!4,colframe=teal!35!black,title={Before you move on}]
What does \te{తొలగించు} mean? (\textquotedblleft{}To remove, to delete\textquotedblright{}.) Say it once more.
\end{tcolorbox}
\section[\texorpdfstring{\te{గుర్తుపట్టు} (gurtupaṭṭu)}{గుర్తుపట్టు (gurtupaṭṭu)}]{\te{గుర్తుపట్టు} (gurtupaṭṭu) — to recognize}

\label{lesson:TE-C283-gurtupattu}



\begin{quote}
Before the new one: say the Telugu for to charge, then the Telugu for to remove.
\end{quote}

\subsection*{You'll want to know: \te{గుర్తుపట్టు}}
\textbf{\te{గుర్తుపట్టు}} — \emph{gurtupaṭṭu} — \textquotedblleft{}to recognize\textquotedblright{}.

\begin{grammarlens}[title={What you've built}]
One more word for this chapter.
\end{grammarlens}

\subsection*{Your turn}
\begin{itemize}
\raggedright
  \item \emph{Say it:} \emph{gurtupaṭṭu}
  \item \emph{Say it:} \emph{gurtupaṭṭu}, once more
  \item \emph{Say it:} say \emph{gurtupaṭṭu}
  \item \emph{Recall:} say the Telugu for to charge, then the Telugu for to remove, then say \emph{gurtupaṭṭu} again
\end{itemize}

\begin{tcolorbox}[breakable,colback=teal!4,colframe=teal!35!black,title={Before you move on}]
What does \te{గుర్తుపట్టు} mean? (\textquotedblleft{}To recognize\textquotedblright{}.) Say it once more.
\end{tcolorbox}
\section[\texorpdfstring{\te{వ్యతిరేకించు} (vyatirēkiñcu)}{వ్యతిరేకించు (vyatirēkiñcu)}]{\te{వ్యతిరేకించు} (vyatirēkiñcu) — to oppose}

\label{lesson:TE-C283-vyatirekincu}



\begin{quote}
Before the new one: say the Telugu for to remove, then the Telugu for to recognize.
\end{quote}

\subsection*{You'll want to know: \te{వ్యతిరేకించు}}
\textbf{\te{వ్యతిరేకించు}} — \emph{vyatirēkiñcu} — \textquotedblleft{}to oppose\textquotedblright{}.

\begin{grammarlens}[title={What you've built}]
One more word for this chapter.
\end{grammarlens}

\subsection*{Your turn}
\begin{itemize}
\raggedright
  \item \emph{Say it:} \emph{vyatirēkiñcu}
  \item \emph{Say it:} \emph{vyatirēkiñcu}, once more
  \item \emph{Say it:} say \emph{vyatirēkiñcu}
  \item \emph{Recall:} say the Telugu for to remove, then the Telugu for to recognize, then say \emph{vyatirēkiñcu} again
\end{itemize}

\begin{tcolorbox}[breakable,colback=teal!4,colframe=teal!35!black,title={Before you move on}]
What does \te{వ్యతిరేకించు} mean? (\textquotedblleft{}To oppose\textquotedblright{}.) Say it once more.
\end{tcolorbox}
\section[\texorpdfstring{\te{వాదించు} (vādiñcu)}{వాదించు (vādiñcu)}]{\te{వాదించు} (vādiñcu) — to argue}

\label{lesson:TE-C283-vadincu}



\begin{quote}
Before the new one: say the Telugu for to recognize, then the Telugu for to oppose.
\end{quote}

\subsection*{You'll want to know: \te{వాదించు}}
\textbf{\te{వాదించు}} — \emph{vādiñcu} — \textquotedblleft{}to argue\textquotedblright{}.

\begin{grammarlens}[title={What you've built}]
One more word for this chapter.
\end{grammarlens}

\subsection*{Your turn}
\begin{itemize}
\raggedright
  \item \emph{Say it:} \emph{vādiñcu}
  \item \emph{Say it:} \emph{vādiñcu}, once more
  \item \emph{Say it:} say \emph{vādiñcu}
  \item \emph{Recall:} say the Telugu for to recognize, then the Telugu for to oppose, then say \emph{vādiñcu} again
\end{itemize}

\begin{tcolorbox}[breakable,colback=teal!4,colframe=teal!35!black,title={Before you move on}]
What does \te{వాదించు} mean? (\textquotedblleft{}To argue\textquotedblright{}.) Say it once more.
\end{tcolorbox}
